\documentclass[12pt, a4paper]{article}

% ── Packages ──────────────────────────────────────────────────────────────────
\usepackage{amsmath, amssymb, amsthm}
\usepackage{geometry}
\usepackage{booktabs}
\usepackage{array}
\usepackage{hyperref}
\usepackage{xcolor}
\usepackage{tcolorbox}
\usepackage{enumitem}
\usepackage{parskip}

\geometry{margin=2.5cm}

% ══════════════════════════════════════════════════════════════════════════════
% ── Student-mode switch ───────────────────────────────────────────────────────
%
%   \studentmodetrue  → student version: blanks replace selected equations
%   \studentmodefalse → instructor version: all equations shown in full
%
\newif\ifstudentmode
\studentmodefalse       % default: instructor mode

\ifx\studentflag\undefined
  % no flag passed — keep studentmodefalse
\else
  \studentmodetrue
\fi
% ══════════════════════════════════════════════════════════════════════════════

% ── Student-mode macros ───────────────────────────────────────────────────────
%
%  \sol[<width>]{<math content>}
%    Use INSIDE a math environment (display or inline).
%    Student mode  → underlined blank of <width> (default 6cm).
%    Instructor mode → renders <math content> normally.
%    Example:  \[ J(w,b) = \sol{\frac{1}{n}\sum(y_i - wx_i - b)^2} \]
%
\newcommand{\sol}[2][6cm]{%
  \ifstudentmode
    \underline{\hspace{#1}}%
  \else
    {#2}%
  \fi
}

%  \soleq[<height>]{<math content>}
%    Replaces an ENTIRE displayed equation with a blank ruled box.
%    Student mode  → two horizontal rules with <height> gap (default 3cm).
%    Instructor mode → renders \[ <math content> \] normally.
%    Example:  \soleq{\bth = (\bX^\top\bX)^{-1}\bX^\top\by}
%
\newcommand{\soleq}[2][3cm]{%
  \ifstudentmode
    \par\vspace{4pt}%
    \noindent\rule{\linewidth}{0.5pt}\\[#1]%
    \noindent\rule{\linewidth}{0.5pt}%
    \par\vspace{4pt}%
  \else
    \[#2\]%
  \fi
}
% ──────────────────────────────────────────────────────────────────────────────

% ── Colors & boxes ────────────────────────────────────────────────────────────
\definecolor{mainblue}{RGB}{30, 100, 180}
\definecolor{boxbg}{RGB}{235, 243, 255}
\definecolor{notebg}{RGB}{255, 248, 230}

\tcbuselibrary{skins, breakable}

\newtcolorbox{resultbox}{
  colback=boxbg, colframe=mainblue,
  boxrule=1.2pt, arc=4pt,
  title={Result},
  fonttitle=\bfseries\color{white},
  coltitle=white, attach boxed title to top left={yshift=-2mm, xshift=4mm},
  boxed title style={colback=mainblue, arc=3pt}
}

\newtcolorbox{notebox}{
  colback=notebg, colframe=orange!70!black,
  boxrule=1pt, arc=4pt,
  left=4pt, right=4pt
}

% ── Theorem environments ───────────────────────────────────────────────────────
\theoremstyle{definition}
\newtheorem{example}{Example}[section]

% ── Math macros ───────────────────────────────────────────────────────────────
\newcommand{\bth}{\boldsymbol{\theta}}
\newcommand{\bx}{\mathbf{x}}
\newcommand{\bX}{\mathbf{X}}
\newcommand{\by}{\mathbf{y}}
\newcommand{\bw}{\mathbf{w}}
\newcommand{\xbar}{\bar{x}}
\newcommand{\ybar}{\bar{y}}
\newcommand{\yhat}{\hat{y}}

% ── Document ──────────────────────────────────────────────────────────────────
\begin{document}

\begin{center}
  {\LARGE \textbf{Section 1.1: Linear Regression}} \\[0.5em]
  {\large AI Class — Lecture Notes}%
  \ifstudentmode
    \quad {\normalsize\color{gray}[Student Copy]}%
  \fi
\end{center}

\vspace{1em}
\hrule
\vspace{1.5em}

% ══════════════════════════════════════════════════════════════════════════════
\section{Problem Setup}

Suppose we observe $n$ data points $(x_i, y_i)$ and believe there is an
approximately linear relationship between $x$ and $y$, corrupted by noise.
For example:

\begin{center}
\begin{tabular}{ccccc}
\toprule
$x$ & 1 & 2 & 3 & 4 \\
\midrule
$y$ & 1 & 3 & 3 & 5 \\
\bottomrule
\end{tabular}
\end{center}

We want to find a line
\[
  \yhat = wx + b
\]
that best fits this data, where $w$ is the \textbf{slope (weight)} and $b$
is the \textbf{intercept (bias)}.

Because of noise, no single line passes through all four points exactly.
We therefore need a way to quantify \emph{how wrong} our predictions are, so
that we can find the \emph{least wrong} line.

% ══════════════════════════════════════════════════════════════════════════════
\section{Cost Function: Mean Squared Error}

For a given choice of $w$ and $b$, the \textbf{residual} at point $i$ is
\[
  e_i = y_i - \yhat_i = y_i - wx_i - b.
\]

We measure the overall fit using the \textbf{Mean Squared Error (MSE)}:

% ── \sol applied: student fills in the MSE formula ───────────────────────────
\begin{equation}
  J(w, b) = \sol[9cm]{\frac{1}{n}\sum_{i=1}^{n}(y_i - wx_i - b)^2}
  \label{eq:mse}
\end{equation}
% ─────────────────────────────────────────────────────────────────────────────

Several properties make MSE a natural choice:
\begin{itemize}[nosep]
  \item Squaring makes all errors positive, so positive and negative residuals
        do not cancel.
  \item Larger errors are penalised more than smaller ones (quadratic penalty).
  \item $J = 0$ is a perfect fit; otherwise the larger $J$, the worse the fit.
  \item $J(w,b)$ is a smooth, bowl-shaped (convex) surface in the
        $(w, b)$ plane, guaranteeing a unique global minimum.
\end{itemize}

\textbf{Goal:} find the values of $w$ and $b$ that minimise $J(w,b)$.

% ══════════════════════════════════════════════════════════════════════════════
\section{Minimising MSE: Scalar Calculus Derivation}

Because $J$ is smooth and convex, the minimum occurs where both partial
derivatives vanish simultaneously:
\[
  \frac{\partial J}{\partial b} = 0
  \qquad \text{and} \qquad
  \frac{\partial J}{\partial w} = 0.
\]
This gives two equations in two unknowns, which we solve sequentially.

% ──────────────────────────────────────────────────────────────────────────────
\subsection{Step 1: Solve for \texorpdfstring{$b$}{b}}

Differentiating \eqref{eq:mse} with respect to $b$:
\[
  \frac{\partial J}{\partial b}
  = \frac{1}{n}\sum_{i=1}^{n} 2(y_i - wx_i - b)(-1)
  = \frac{-2}{n}\sum_{i=1}^{n}(y_i - wx_i - b) = 0.
\]

Expanding and dividing by $n$:
\[
  \frac{1}{n}\sum_{i=1}^{n} y_i
  = w \cdot \frac{1}{n}\sum_{i=1}^{n} x_i + b
  \quad \Longrightarrow \quad
  \ybar = w\xbar + b.
\]

% ── \sol applied: student fills in the result for b ──────────────────────────
\begin{resultbox}
\[
  b = \sol[7cm]{\ybar - w\xbar}
\]
\end{resultbox}
% ─────────────────────────────────────────────────────────────────────────────

\begin{notebox}
\textbf{Interpretation.} For any slope $w$, the optimal line always passes
through the \emph{centroid} $(\xbar, \ybar)$ of the dataset.
\end{notebox}

% ──────────────────────────────────────────────────────────────────────────────
\subsection{Step 2: Solve for \texorpdfstring{$w$}{w}}

Differentiating \eqref{eq:mse} with respect to $w$:
\[
  \frac{\partial J}{\partial w}
  = \frac{-2}{n}\sum_{i=1}^{n}(y_i - wx_i - b)\,x_i = 0.
\]

Substituting $b = \ybar - w\xbar$ and letting $\tilde{x}_i = x_i - \xbar$,
$\tilde{y}_i = y_i - \ybar$:

\[
  \sum_{i=1}^{n}(y_i - wx_i - \ybar + w\xbar)\,x_i = 0
  \quad\Longrightarrow\quad
  \sum_{i=1}^{n}\bigl(\tilde{y}_i - w\tilde{x}_i\bigr)\,x_i = 0.
\]

Expanding and noting that $\sum_i \tilde{x}_i = 0$:

\[
  \sum_{i=1}^{n}\tilde{y}_i\, x_i = w\sum_{i=1}^{n}\tilde{x}_i\, x_i
  \quad\Longrightarrow\quad
  \sum_{i=1}^{n}\tilde{x}_i\tilde{y}_i = w\sum_{i=1}^{n}\tilde{x}_i^2.
\]

% ── \sol applied: student fills in the result for w ──────────────────────────
\begin{resultbox}
\[
  w = \sol[8cm]{\frac{\displaystyle\sum_{i=1}^{n}(x_i - \xbar)(y_i - \ybar)}
                     {\displaystyle\sum_{i=1}^{n}(x_i - \xbar)^2}}
\]
\end{resultbox}
% ─────────────────────────────────────────────────────────────────────────────

\begin{notebox}
\textbf{Interpretation.} The optimal slope equals the \emph{sample covariance}
of $x$ and $y$ divided by the \emph{sample variance} of $x$.
A positive covariance (both tend to increase together) yields a positive slope,
and vice versa.
\end{notebox}

% ──────────────────────────────────────────────────────────────────────────────
\subsection{Numerical Result}

For the dataset above, $\xbar = 2.5$ and $\ybar = 3$:

\begin{center}
\renewcommand{\arraystretch}{1.3}
\begin{tabular}{ccccc}
\toprule
$i$ & $x_i - \xbar$ & $y_i - \ybar$
    & $(x_i-\xbar)(y_i-\ybar)$ & $(x_i-\xbar)^2$ \\
\midrule
1 & $-1.5$ & $-2$ & $3.0$ & $2.25$ \\
2 & $-0.5$ & $\phantom{-}0$ & $0.0$ & $0.25$ \\
3 & $\phantom{-}0.5$ & $\phantom{-}0$ & $0.0$ & $0.25$ \\
4 & $\phantom{-}1.5$ & $\phantom{-}2$ & $3.0$ & $2.25$ \\
\midrule
  & & \textbf{Sum} & $6$ & $5$ \\
\bottomrule
\end{tabular}
\end{center}

% ── \sol applied: student computes w and b from the sums ─────────────────────
\[
  w = \sol{\dfrac{6}{5} = 1.2},
  \qquad
  b = \sol{3 - 1.2 \times 2.5 = 0}.
\]

The fitted line is $\yhat = \sol{1.2\,x}$, with:

\[
  \mathrm{MSE}
  = \frac{(-0.2)^2 + (0.6)^2 + (-0.6)^2 + (0.2)^2}{4}
  = \sol{\dfrac{0.8}{4} = 0.2} \neq 0.
\]
% ─────────────────────────────────────────────────────────────────────────────

The non-zero MSE reflects irreducible noise in the data — not a failure of the
method.

% ══════════════════════════════════════════════════════════════════════════════
\section{Generalisation: Multiple Features}

With a single feature we derived two scalar equations. In practice, a data
point often has $d$ features $x_1, x_2, \ldots, x_d$, and the model becomes:
\[
  \yhat = w_1 x_1 + w_2 x_2 + \cdots + w_d x_d + b.
\]

Taking a partial derivative for each of the $d+1$ parameters and setting them
to zero would produce $d+1$ scalar equations with the same structure. Rather
than repeating this process, we rewrite everything in \textbf{matrix form},
which handles all parameters simultaneously.

% ──────────────────────────────────────────────────────────────────────────────
\subsection{Matrix Notation}

Collect all parameters into a single vector:
\[
  \bth = \begin{bmatrix} b \\ w_1 \\ \vdots \\ w_d \end{bmatrix}
  \in \mathbb{R}^{d+1}.
\]

Stack all $n$ data points row-by-row into the \textbf{design matrix}
$\bX \in \mathbb{R}^{n \times (d+1)}$, prepending a column of ones to account
for the bias:
\[
  \bX =
  \begin{bmatrix}
    1 & x_1^{(1)} & x_2^{(1)} & \cdots & x_d^{(1)} \\
    1 & x_1^{(2)} & x_2^{(2)} & \cdots & x_d^{(2)} \\
    \vdots & \vdots & \vdots & \ddots & \vdots \\
    1 & x_1^{(n)} & x_2^{(n)} & \cdots & x_d^{(n)}
  \end{bmatrix},
  \qquad
  \by =
  \begin{bmatrix} y^{(1)} \\ y^{(2)} \\ \vdots \\ y^{(n)} \end{bmatrix}.
\]

With this notation, the vector of all predictions is simply $\bX\bth$, and the
MSE cost becomes:

% ── \sol applied: student writes the matrix form of MSE ──────────────────────
\[
  J(\bth) = \sol[8cm]{\frac{1}{n}\left\|\by - \bX\bth\right\|^2}
\]
% ─────────────────────────────────────────────────────────────────────────────

% ──────────────────────────────────────────────────────────────────────────────
\subsection{The Normal Equations}

Setting the gradient of $J$ with respect to $\bth$ to zero and expanding
step by step:

\begin{align*}
  &\nabla_{\bth}\, J = \mathbf{0} \\[6pt]
  \Rightarrow\;&
    \nabla_{\bth}\,\frac{1}{n}\|\bX\bth - \by\|_2^2 = \mathbf{0} \\[6pt]
  \Rightarrow\;&
    \nabla_{\bth}
    \left(\bX\bth - \by\right)^{\!\top}
    \left(\bX\bth - \by\right) = \mathbf{0} \\[6pt]
  \Rightarrow\;&
% ── \sol applied: student expands the quadratic form ────────────────────────
    \nabla_{\bth}\!\left(
      \sol[9cm]{\bth^{\top}\bX^{\top}\bX\bth
      - 2\bth^{\top}\bX^{\top}\by
      + \by^{\top}\by}
    \right) = \mathbf{0} \\[6pt]
% ────────────────────────────────────────────────────────────────────────────
  \Rightarrow\;&
% ── \sol applied: student applies the gradient ───────────────────────────────
    \sol[7cm]{2\bX^{\top}\bX\bth - 2\bX^{\top}\by} = \mathbf{0} \\[6pt]
% ────────────────────────────────────────────────────────────────────────────
  \Rightarrow\;&
    \bX^{\top}\bX\,\bth = \bX^{\top}\by.
\end{align*}

Provided $\bX^{\top}\bX$ is invertible, the unique solution is:

% ── \sol applied: student writes the normal equations result ─────────────────
\begin{resultbox}
\[
  \bth = \sol[8cm]{\left(\bX^{\top}\bX\right)^{-1}\bX^{\top}\by}
\]
\end{resultbox}
% ─────────────────────────────────────────────────────────────────────────────

This is the \textbf{Normal Equation} — a single closed-form formula that solves
linear regression for any number of features $d$ and any number of data points
$n$ (provided $n > d+1$ and the columns of $\bX$ are linearly independent).

% ══════════════════════════════════════════════════════════════════════════════
\section{Example: Two Binary Features}

\begin{example}
Consider a dataset with two binary inputs where the true relationship is
$y = x_1 + x_2$:

\begin{center}
\begin{tabular}{ccc}
\toprule
$x_1$ & $x_2$ & $y$ \\
\midrule
0 & 0 & 0 \\
0 & 1 & 1 \\
1 & 0 & 1 \\
1 & 1 & 2 \\
\bottomrule
\end{tabular}
\end{center}

The design matrix and target vector are:
\[
  \bX =
  \begin{bmatrix} 1 & 0 & 0 \\ 1 & 0 & 1 \\ 1 & 1 & 0 \\ 1 & 1 & 1 \end{bmatrix},
  \qquad
  \by =
  \begin{bmatrix} 0 \\ 1 \\ 1 \\ 2 \end{bmatrix}.
\]

Computing:
\[
  \bX^{\top}\bX =
  \begin{bmatrix} 4 & 2 & 2 \\ 2 & 2 & 1 \\ 2 & 1 & 2 \end{bmatrix},
  \qquad
  \bX^{\top}\by =
  \begin{bmatrix} 4 \\ 3 \\ 3 \end{bmatrix}.
\]

Solving $(\bX^{\top}\bX)\,\bth = \bX^{\top}\by$ gives:

% ── \soleq applied: student writes the final θ solution ──────────────────────
\soleq[2.5cm]{%
  \bth =
  \begin{bmatrix} b \\ w_1 \\ w_2 \end{bmatrix}
  =
  \begin{bmatrix} 0 \\ 1 \\ 1 \end{bmatrix}
}
% ─────────────────────────────────────────────────────────────────────────────

so $\yhat = \sol{x_1 + x_2}$ with $\mathrm{MSE} = \sol[2cm]{0}$.

The normal equations recover the exact underlying relationship because the
data is truly linear.
\end{example}

\begin{notebox}
\textbf{Looking ahead.}
In Section~1.3 we will use the same four input points
$\{(0,0),\,(0,1),\,(1,0),\,(1,1)\}$ but change the output labels to
$\{0,\,1,\,1,\,0\}$ (the XOR function). Applying the normal equations to
that dataset yields $w_1 = w_2 = 0$, $b = \tfrac{1}{2}$ — the model simply
predicts the mean $0.5$ everywhere, achieving $\mathrm{MSE} = 0.25$.
This is the best any linear model can do on XOR, which motivates the
need for non-linearity.
\end{notebox}

% ══════════════════════════════════════════════════════════════════════════════
\section{Summary}

\begin{center}
\renewcommand{\arraystretch}{1.5}
\begin{tabular}{ll}
\toprule
\textbf{Model}             & $\yhat = \bth^{\top}\bx$ \\
\textbf{Loss function}     & MSE: $J = \dfrac{1}{n}\|\by - \bX\bth\|^2$ \\
\textbf{Closed-form solution} & Normal equations: $\bth = (\bX^{\top}\bX)^{-1}\bX^{\top}\by$ \\
\textbf{Works well when}   & The true relationship is (approximately) linear \\
\bottomrule
\end{tabular}
\end{center}

\bigskip
\textbf{Up next (Section~1.2):} What if the output $y$ is a \emph{category}
rather than a continuous number? Predicting a probability requires a
different model and a different loss function — logistic regression with
cross-entropy.

\end{document}
